\documentclass[../main.tex]{subfiles}
\begin{document}
%=========================================================
% TĐ_PBT_IQA_HK1_Tuan17_SoSanhCacSoTrongPhamVi100
% Phiếu bài tập Toán tư duy cấp độ A
% Chủ đề SO SÁNH CÁC SỐ TRONG PHẠM VI 100 (Buổi 1, Buổi 2)
% Soạn theo mẫu (\hktuan, bt, listEX)
% Toàn bộ hình minh họa được vẽ bằng TikZ (không dùng file ảnh), tiền tố lệnh: SS
% Cần các thư viện TikZ có sẵn trong main: shapes.geometric, arrows, calc, patterns
%=========================================================

\usetikzlibrary{shadings}
% Không kéo giãn khoảng cách dọc để dồn trang đầy (khôi phục ở cuối file)
\makeatletter\let\SSluutextbottom\@textbottom\raggedbottom\makeatother

%---------- Lệnh dùng chung cho phiếu (giống các tuần trước) ----------
% Chỗ chấm điền
\providecommand{\TTcham}[1][3cm]{\makebox[#1]{\dotfill}}
% Dòng thông tin học sinh
\providecommand{\TTthongtin}{%
  \noindent\begin{tabular}{|p{7.6cm}|p{4cm}|p{4cm}|}
    \hline
    \textbf{Họ và Tên HS:}\dotfill & \textbf{Lớp:}\dotfill & \textbf{Ngày:}\dotfill \\
    \hline
  \end{tabular}\par\medskip}
% Dòng chủ đề: một dòng, tự thu nhỏ vừa khổ nếu quá dài
\providecommand{\SStieude}{}\renewcommand{\SStieude}[1]{\sbox0{\color{mauchinh!120}\bfseries\Large #1}\ifdim\wd0>\linewidth\resizebox{\linewidth}{!}{\usebox0}\else\usebox0\fi}
% Đánh lại số bài từ 1 cho buổi 2 (chỉ khi môi trường bt có bộ đếm tên bt)
\makeatletter
\providecommand{\TTdatlaibai}{\@ifundefined{c@bt}{}{\setcounter{bt}{0}}}
\makeatother

%---------- Màu dùng cho hình vẽ ----------
\definecolor{SSxanh}{RGB}{100,170,235}\definecolor{SSxanhdam}{RGB}{30,100,180}
\definecolor{SSla}{RGB}{150,205,120}\definecolor{SSladam}{RGB}{45,125,60}
\definecolor{SSvang}{RGB}{255,215,90}\definecolor{SScam}{RGB}{250,150,60}
\definecolor{SSdo}{RGB}{225,70,70}\definecolor{SShong}{RGB}{250,170,200}
\definecolor{SStim}{RGB}{170,140,220}\definecolor{SSnau}{RGB}{125,80,45}
\definecolor{SSkem}{RGB}{255,247,220}\definecolor{SSnuoc}{RGB}{205,232,250}
\definecolor{SSxam}{RGB}{228,232,238}\definecolor{SSnaudam}{RGB}{85,50,25}

%---------- Kiểu dùng chung ----------
\tikzset{
  % thẻ số tròn góc, màu pastel: SSthe=<tên màu>
  SSthe/.style={draw=#1!65!black,thick,fill=#1!25,rounded corners=6pt,minimum width=1.1cm,minimum height=0.95cm,inner sep=0pt,font=\normalfont\bfseries\Large},
  % ô trống để học sinh điền
  SSo/.style={draw,thick,rounded corners=3pt,minimum width=0.95cm,minimum height=0.8cm,inner sep=0pt,fill=white},
}
% Ô trống nằm trong dòng chữ
\providecommand{\SSo}{}\renewcommand{\SSo}{\hspace{4pt}\tikz[baseline=-0.75ex]\node[SSo]{};\hspace{4pt}}
% Ô điền sẵn (mẫu), chữ đỏ: \SSomau{46}
\providecommand{\SSomau}{}\renewcommand{\SSomau}[1]{\hspace{4pt}\tikz[baseline=-0.75ex]\node[SSo]{\textcolor{red}{\bfseries #1}};\hspace{4pt}}
% Hai ý trên một dòng, thẳng cột: \SShaicot{ý trái}{ý phải}
\providecommand{\SShaicot}{}\renewcommand{\SShaicot}[2]{\parbox[t]{0.5\linewidth}{#1}\parbox[t]{0.48\linewidth}{#2}}
% Tên màu pastel xoay vòng theo số thứ tự 0..4
\providecommand{\SSmaule}{}\renewcommand{\SSmaule}[1]{\ifcase#1 SSxanh\or SSla\or SSvang\or SShong\or SStim\fi}

%================ HÌNH VẼ: BUỔI 1 ================
% Khối chục - đơn vị: \SSkhoi{số chục}{số đơn vị} (vẽ từ gốc tọa độ, cao 1,9cm)
\providecommand{\SSkhoi}{}\renewcommand{\SSkhoi}[2]{%
  \ifnum#1>0 \foreach \c in {1,...,#1}{%
    \begin{scope}[shift={(0.28*\c-0.28,0)}]
      \shade[left color=SSxanh!70!white,right color=SSxanh] (0,0) rectangle (0.22,1.9);
      \foreach \r in {1,...,9}{\draw[SSxanhdam!70,very thin] (0,0.19*\r)--(0.22,0.19*\r);}
      \draw[SSxanhdam,thin] (0,0) rectangle (0.22,1.9);
    \end{scope}}\fi
  \ifnum#2>0 \foreach \u in {1,...,#2}{%
    \pgfmathtruncatemacro{\cot}{mod(\u-1,3)}\pgfmathtruncatemacro{\hang}{div(\u-1,3)}%
    \shade[left color=SScam!60!white,right color=SScam] (0.28*#1+0.12+0.25*\cot,0.25*\hang) rectangle ++(0.21,0.21);
    \draw[SScam!45!black,thin] (0.28*#1+0.12+0.25*\cot,0.25*\hang) rectangle ++(0.21,0.21);}\fi}
% Thẻ so sánh bằng hình: \SSsosanh{chục 1}{đv 1}{chục 2}{đv 2}
\providecommand{\SSsosanh}{}\renewcommand{\SSsosanh}[4]{%
  \begin{tikzpicture}[font=\normalfont]
    \draw[rounded corners=8pt,thick,draw=SSxanh!70,fill=SSxanh!8] (-0.3,-1.3) rectangle (7.0,2.3);
    \begin{scope}[shift={(0.1,0)}]\SSkhoi{#1}{#2}\end{scope}
    \begin{scope}[shift={(4.5,0)}]\SSkhoi{#3}{#4}\end{scope}
    \node[SSo,minimum width=1.1cm,minimum height=0.95cm] at (3.4,0.95) {};
    \node[anchor=east,font=\normalfont\small] at (0.45,-0.75) {Số:};\node[SSo] at (1.15,-0.75) {};
    \node[anchor=east,font=\normalfont\small] at (4.85,-0.75) {Số:};\node[SSo] at (5.55,-0.75) {};
  \end{tikzpicture}}

% Cá sấu há miệng (quay mặt sang phải, dài khoảng 4cm, bản lề tại gốc tọa độ)
\providecommand{\SScasau}{}\renewcommand{\SScasau}{%
  \begin{scope}[line join=round]
    % thân và đuôi phía sau đầu
    \fill[SSla!80!SSladam,draw=SSladam!80!black,thick] (-0.1,0.5) .. controls (-0.8,0.52) and (-1.5,0.25) .. (-2.4,-0.1) .. controls (-1.5,-0.3) and (-0.8,-0.38) .. (-0.1,-0.5) -- cycle;
    % gai lưng
    \foreach \x in {-0.35,-0.75,-1.15,-1.55,-1.95}{\fill[SSladam!85!black] (\x-0.1,0.4+0.1*\x)--(\x,0.56+0.1*\x)--(\x+0.1,0.4+0.1*\x)--cycle;}
    % hàm dưới
    \fill[SSla!75!SSladam,draw=SSladam!80!black,thick] (-0.3,0.02) .. controls (-0.4,-0.4) and (0.2,-0.55) .. (0.8,-0.5) .. controls (1.2,-0.47) and (1.5,-0.38) .. (1.75,-0.3) .. controls (1.95,-0.25) and (1.95,0.0) .. (1.75,0.02) -- cycle;
    \fill[pink!80!red!60] (0.1,0.0) .. controls (0.5,-0.2) and (1.2,-0.18) .. (1.6,-0.04) -- (1.6,0.02) -- cycle;
    \foreach \x in {0.55,0.85,1.15,1.45}{\fill[white,draw=gray!70,very thin] (\x-0.07,0.02)--(\x,0.16)--(\x+0.07,0.02)--cycle;}
    % hàm trên mở ra 30 độ
    \begin{scope}[rotate around={30:(0,0)}]
      \foreach \x in {0.55,0.85,1.15,1.45}{\fill[white,draw=gray!70,very thin] (\x-0.07,0.0)--(\x,-0.14)--(\x+0.07,0.0)--cycle;}
      \fill[SSla,draw=SSladam!80!black,thick] (0,0) .. controls (-0.3,0.05) and (-0.3,0.6) .. (0.1,0.65) .. controls (0.5,0.75) and (1.0,0.5) .. (1.6,0.38) .. controls (1.95,0.34) and (1.95,0.04) .. (1.7,0.0) -- cycle;
      \fill[white] (0.3,0.5) circle (0.15);\fill[black] (0.34,0.5) circle (0.08);\fill[white] (0.37,0.54) circle (0.025);
      \fill[SSladam] (1.7,0.24) circle (0.035);
      \draw[SSladam!70!black,thin,line cap=round] (0.12,0.66)--(0.0,0.72) (0.3,0.7)--(0.3,0.78);
    \end{scope}
  \end{scope}}
% Bong bóng số lớn / bé dùng trong khung mẹo: \SSbongso{số}{màu}
\providecommand{\SSbongso}{}\renewcommand{\SSbongso}[2]{%
  \node[SSthe=#2,circle,minimum size=1.15cm] {#1};}

% Một hàng số để khoanh: \SShangso{a,b,c,...}
\providecommand{\SShangso}{}\renewcommand{\SShangso}[1]{%
  \begin{tikzpicture}
    \foreach \x [count=\k] in {#1}{%
      \pgfmathtruncatemacro{\m}{mod(\k-1,5)}\edef\mc{\SSmaule{\m}}%
      \node[SSthe=\mc] at (0.55+1.28*\k-1.28,0) {\x};}
  \end{tikzpicture}}

% Hàng số có một số đã khoanh sẵn (mẫu): \SShangsomau{a,b,c,...}{thứ tự số được khoanh}
\providecommand{\SShangsomau}{}\renewcommand{\SShangsomau}[2]{%
  \begin{tikzpicture}
    \useasboundingbox (0,-0.475) rectangle (6.9,0.475);% chữ "Mẫu" nhô lên, không làm hàng cao hơn
    \foreach \x [count=\k] in {#1}{%
      \pgfmathtruncatemacro{\m}{mod(\k-1,5)}\edef\mc{\SSmaule{\m}}%
      \node[SSthe=\mc] at (0.55+1.28*\k-1.28,0) {\x};}
    \draw[red,very thick] (0.55+1.28*#2-1.28,0) ellipse (0.78 and 0.68);
    \node[orange,font=\normalfont\bfseries\normalsize] at (0.55+1.28*#2-1.28,0.82) {Mẫu};
  \end{tikzpicture}}

% Sâu đo: \SSsaudohang{màu}{5 số}{chữ ô đầu}{chữ ô cuối}{1 = ô to dần, 0 = ô nhỏ dần}
\providecommand{\SSsaudohang}{}\renewcommand{\SSsaudohang}[5]{%
  \begin{tikzpicture}[font=\normalfont\bfseries\large,scale=0.92,transform shape]
    % chân
    \foreach \k in {1,...,5}{\draw[#1!70!black,line width=1.5pt,line cap=round] (0.95*\k+0.05,-0.4)--(0.95*\k-0.08,-0.72) (0.95*\k+0.4,-0.4)--(0.95*\k+0.5,-0.7);}
    % thân (màu pastel)
    \foreach \k in {1,...,5}{\shade[ball color=#1!40!white] (0.95*\k+0.2,0) circle (0.52);\draw[#1!70!white,thin] (0.95*\k+0.2,0) circle (0.52);}
    \foreach \x [count=\k] in {#2}{\node at (0.95*\k+0.2,0) {\x};}
    % râu
    \draw[#1!75!black,thick,line cap=round] (-0.12,0.55) .. controls (-0.25,0.95) and (-0.05,1.1) .. (0.12,1.05);
    \draw[#1!75!black,thick,line cap=round] (0.2,0.6) .. controls (0.2,0.95) and (0.4,1.1) .. (0.55,1.02);
    \fill[#1!75!black] (0.12,1.05) circle (0.06) (0.55,1.02) circle (0.06);
    % đầu
    \shade[ball color=#1!75!white] (0,0) circle (0.66);
    \draw[#1!70!white,thin] (0,0) circle (0.66);
    \fill[white] (-0.12,0.12) circle (0.17);\fill[black] (-0.16,0.12) circle (0.085);\fill[white] (-0.19,0.16) circle (0.025);
    \fill[SShong,opacity=0.7] (-0.38,-0.17) ellipse (0.12 and 0.07);
    \draw[#1!75!black,thick,line cap=round] (-0.5,-0.05) .. controls (-0.38,-0.34) and (-0.05,-0.36) .. (0.08,-0.15);
    % ô trống: to dần (bé -> lớn) hoặc nhỏ dần (lớn -> bé)
    \ifnum#5=1 \def\SSds{0.7/6.6,0.85/7.38,1.0/8.31,1.15/9.39,1.3/10.62}\def\SSxa{6.95}\def\SSxb{11.27}%
    \else \def\SSds{1.3/6.6,1.15/7.98,1.0/9.21,0.85/10.29,0.7/11.22}\def\SSxa{7.25}\def\SSxb{11.57}\fi
    \foreach \sz/\xo in \SSds {\draw[thick,rounded corners=3pt,fill=white] (\xo,-\sz/2) rectangle ++(\sz,\sz);}
    % (không ghi chữ "bé nhất / lớn nhất" dưới ô; #3, #4 giữ lại để tương thích lời gọi)
    \draw[-latex,thick,gray] (6.6,0.85)--(12.1,0.85);
  \end{tikzpicture}}

% ---- Cửa hàng trái cây ----
\providecommand{\SStaoicon}{}\renewcommand{\SStaoicon}{%
  \shade[ball color=red!80!orange] (0,0) circle (0.3);
  \draw[SSnau,thick] (0,0.27)--(0.03,0.4);
  \fill[SSla!80!SSladam,rotate around={30:(0.12,0.38)}] (0.12,0.38) ellipse (0.11 and 0.05);
  \fill[white,opacity=0.5] (-0.13,0.1) ellipse (0.04 and 0.08);}
\providecommand{\SScamicon}{}\renewcommand{\SScamicon}{%
  \shade[ball color=orange] (0,0) circle (0.3);
  \fill[SSla!80!SSladam,rotate around={20:(0.0,0.3)}] (0,0.3) ellipse (0.1 and 0.04);
  \fill[orange!50!brown] (0.1,0.0) circle (0.012) (-0.05,-0.12) circle (0.012) (-0.12,0.05) circle (0.012) (0.0,0.13) circle (0.012) (0.14,-0.1) circle (0.012);
  \fill[white,opacity=0.5] (-0.13,0.1) ellipse (0.04 and 0.08);}
\providecommand{\SSleicon}{}\renewcommand{\SSleicon}{%
  \shade[ball color=yellow!65!SSla] (0,0.28) .. controls (-0.08,0.22) and (-0.25,0.05) .. (-0.27,-0.1) .. controls (-0.29,-0.35) and (0.29,-0.35) .. (0.27,-0.1) .. controls (0.25,0.05) and (0.08,0.22) .. (0,0.28);
  \draw[SSnau,thick] (0,0.27)--(0.04,0.4);
  \fill[SSla!80!SSladam,rotate around={25:(0.1,0.37)}] (0.1,0.37) ellipse (0.1 and 0.04);
  \fill[white,opacity=0.45] (-0.13,-0.05) ellipse (0.035 and 0.08);}
\providecommand{\SSdauicon}{}\renewcommand{\SSdauicon}{%
  \shade[ball color=red!85!black] (0,-0.28) .. controls (-0.3,-0.05) and (-0.28,0.2) .. (-0.12,0.2) .. controls (-0.05,0.22) and (0.05,0.22) .. (0.12,0.2) .. controls (0.28,0.2) and (0.3,-0.05) .. (0,-0.28);
  \foreach \x/\y in {-0.1/0.08,0.0/0.1,0.1/0.08,-0.14/-0.03,-0.04/-0.01,0.06/-0.01,0.15/-0.03,-0.08/-0.12,0.03/-0.12,0.1/-0.13,-0.02/-0.2}{\fill[yellow!70!orange] (\x,\y) ellipse (0.016 and 0.024);}
  \fill[SSladam] (0,0.2) -- (-0.2,0.25) -- (-0.06,0.27) -- (-0.1,0.36) -- (0,0.29) -- (0.1,0.36) -- (0.06,0.27) -- (0.2,0.25) -- cycle;}
% Thùng gỗ đựng trái cây: \SSthung{biểu tượng}{tên}{số}
\providecommand{\SSthung}{}\renewcommand{\SSthung}[3]{%
  \begin{tikzpicture}[font=\normalfont]
    \fill[black,opacity=0.12] (1.6,-0.08) ellipse (1.75 and 0.14);
    % đống quả
    \foreach \x/\y in {0.55/0.95,1.25/0.95,1.95/0.95,2.65/0.95,0.9/1.45,1.6/1.45,2.3/1.45,1.25/1.95,1.95/1.95}{\begin{scope}[shift={(\x,\y)}]#1\end{scope}}
    % thùng
    \shade[top color=SSnau!55!SSvang,bottom color=SSnau!70!white] (0,0) rectangle (3.2,1.1);
    \draw[SSnaudam,thick] (0,0) rectangle (3.2,1.1);
    \draw[SSnaudam,thin] (0,0.37)--(3.2,0.37) (0,0.73)--(3.2,0.73);
    \foreach \x in {0.12,3.08}{\fill[SSnaudam] (\x,0.2) circle (0.03) (\x,0.55) circle (0.03) (\x,0.9) circle (0.03);}
    \node[draw=SSnaudam,fill=white,rounded corners=3pt,minimum width=2.7cm,minimum height=0.6cm,inner sep=1pt,font=\normalfont\bfseries\footnotesize] at (1.6,0.55) {#2: #3 quả};
  \end{tikzpicture}}

% Bục nhận giải: \SSbuc  (3 bục: hạng 2 - hạng 1 - hạng 3), có ô viết tên và điểm
\providecommand{\SSbuc}{}\renewcommand{\SSbuc}{%
  \begin{tikzpicture}[font=\normalfont\bfseries]
    % ô viết tên
    \foreach \x/\y in {2.1/2.3,5.3/3.2,8.5/1.9}{
      \node[draw,thick,rounded corners=4pt,fill=white,minimum width=2.6cm,minimum height=0.75cm] at (\x,\y+1.0) {};}
    % bục
    \shade[top color=gray!15,bottom color=gray!55] (0.8,0) rectangle (3.4,1.7);
    \shade[top color=yellow!35,bottom color=orange!70!yellow] (3.9,0) rectangle (6.7,2.6);
    \shade[top color=orange!25,bottom color=brown!60] (7.2,0) rectangle (9.8,1.3);
    \draw[thick,draw=black!70] (0.8,0) rectangle (3.4,1.7) (3.9,0) rectangle (6.7,2.6) (7.2,0) rectangle (9.8,1.3);
    \node[font=\normalfont\bfseries\Huge,text=black!55] at (2.1,0.85) {2};
    \node[font=\normalfont\bfseries\Huge,text=orange!60!brown] at (5.3,1.3) {1};
    \node[font=\normalfont\bfseries\Huge,text=brown!70!black] at (8.5,0.65) {3};
    % huy chương
    \foreach \x/\y/\c in {2.1/2.0/gray!55,5.3/2.9/yellow!80!orange,8.5/1.6/brown!55!orange}{
      \begin{scope}[shift={(\x,\y-0.05)}]
        \fill[red!70,draw=red!50!black,thin] (-0.14,0.18)--(0,-0.04)--(0.14,0.18)--cycle;
        \shade[ball color=\c] (0,-0.1) circle (0.14);
      \end{scope}}
    % cờ chiến thắng trên bục hạng nhất
    \draw[thick,SSnau] (6.35,2.6)--(6.35,3.7);\fill[SSdo,draw=red!50!black,thin] (6.35,3.7)--(7.0,3.5)--(6.35,3.3)--cycle;
  \end{tikzpicture}}

%================ HÌNH VẼ: BUỔI 2 ================
% Bọ rùa nhìn từ trên xuống, đầu hướng lên (vẽ tại gốc tọa độ, cao khoảng 0,8cm)
\providecommand{\SSbochua}{}\renewcommand{\SSbochua}{%
  \fill[black] (0,0.34) circle (0.16);
  \draw[black,thick,line cap=round] (-0.05,0.46)--(-0.15,0.6) (0.05,0.46)--(0.15,0.6);
  \shade[ball color=red!85!black] (0,0) ellipse (0.3 and 0.33);
  \draw[black,thick] (0,-0.32)--(0,0.31);
  \fill[black] (-0.15,0.1) circle (0.045) (0.15,0.1) circle (0.045) (-0.17,-0.1) circle (0.04) (0.17,-0.1) circle (0.04) (-0.07,-0.22) circle (0.03) (0.07,-0.22) circle (0.03);
  \fill[white,opacity=0.5] (-0.2,0.2) ellipse (0.04 and 0.07);}
% Ngôi sao / huy hiệu có số (số lẻ), bóng bay có số (số chẵn)
\providecommand{\SSbongbay}{}\renewcommand{\SSbongbay}[3]{% #1 = vị trí, #2 = màu, #3 = số
  \begin{scope}[shift={(#1)}]
    \draw[#2!50!black,thin] (0,-0.62) .. controls (-0.12,-0.8) and (0.12,-0.9) .. (0,-1.1);
    \fill[#2!65!black] (-0.07,-0.7)--(0.07,-0.7)--(0,-0.58)--cycle;
    \shade[shading=radial,inner color=#2!12!white,outer color=#2!70!white] (0,0) ellipse (0.55 and 0.68);
    \draw[#2!60!black,thin] (0,0) ellipse (0.55 and 0.68);
    \fill[white,opacity=0.6] (-0.25,0.3) ellipse (0.08 and 0.14);
    \node[font=\normalfont\bfseries\Large] at (0,0) {#3};
  \end{scope}}
\providecommand{\SShuyhieu}{}\renewcommand{\SShuyhieu}[3]{%
  \node[star,star points=12,star point ratio=1.16,draw=#2!60!black,thick,fill=#2!35,minimum size=1.35cm,inner sep=0pt,font=\normalfont\bfseries\Large] at (#1) {#3};}
% Hàng bóng bay (số chẵn): \SSbonghang{a,b,...}; hàng huy hiệu (số lẻ): \SShuyhieuhang{a,b,...}
\providecommand{\SSbonghang}{}\renewcommand{\SSbonghang}[1]{%
  \begin{tikzpicture}
    \foreach \x [count=\k] in {#1}{%
      \pgfmathtruncatemacro{\m}{mod(\k-1,5)}\edef\mc{\SSmaule{\m}}%
      \pgfmathtruncatemacro{\cot}{mod(\k-1,8)}\pgfmathtruncatemacro{\hang}{div(\k-1,8)}%
      \SSbongbay{2.05*\cot,-1.85*\hang}{\mc}{\x}}
  \end{tikzpicture}}
\providecommand{\SShuyhieuhang}{}\renewcommand{\SShuyhieuhang}[1]{%
  \begin{tikzpicture}
    \foreach \x [count=\k] in {#1}{%
      \pgfmathtruncatemacro{\m}{mod(\k-1,5)}\edef\mc{\SSmaule{\m}}%
      \pgfmathtruncatemacro{\cot}{mod(\k-1,8)}\pgfmathtruncatemacro{\hang}{div(\k-1,8)}%
      \SShuyhieu{2.05*\cot,-1.6*\hang}{\mc}{\x}}
  \end{tikzpicture}}

% Ô vuông nhỏ để đánh dấu (chẵn / lẻ)
\providecommand{\SSocheck}{}\renewcommand{\SSocheck}{\raisebox{-0.6mm}{\framebox[4mm]{\rule{0pt}{4mm}}}}
% Ghép đôi bọ rùa (phần lõi, không có khung tikzpicture): \SSdoicore{số con}{1 = vẽ sẵn vòng khoanh}
\providecommand{\SSdoicore}{}\renewcommand{\SSdoicore}[2]{%
  \foreach \k in {1,...,#1}{\pgfmathtruncatemacro{\cot}{div(\k-1,2)}\pgfmathtruncatemacro{\hang}{mod(\k-1,2)}%
    \begin{scope}[shift={(1.15*\cot,-1.0*\hang)}]\SSbochua\end{scope}}
  \ifnum#2=1
    \pgfmathtruncatemacro{\sodoi}{div(#1,2)}%
    \ifnum\sodoi>0 \foreach \c in {1,...,\sodoi}{\draw[SSxanhdam,thick,dashed,rounded corners=9pt] (1.15*\c-1.15-0.45,0.68) rectangle (1.15*\c-1.15+0.45,-1.42);}\fi
    \pgfmathtruncatemacro{\ledu}{mod(#1,2)}%
    \ifnum\ledu=1 \pgfmathtruncatemacro{\vitri}{div(#1,2)}\draw[SScam!90!SSdo,very thick,dashed] (1.15*\vitri,0.1) circle (0.62);\fi
  \fi}
% Một ô bài tập ghép đôi: \SSdoibai{số con}
\providecommand{\SSdoibai}{}\renewcommand{\SSdoibai}[1]{%
  \begin{tikzpicture}[font=\normalfont]
    \draw[rounded corners=8pt,thick,draw=SSxanh!70,fill=SSxanh!6] (-0.7,-2.45) rectangle (7.45,1.0);
    \begin{scope}[shift={(0.2,0)}]\SSdoicore{#1}{0}\end{scope}
    \node[anchor=west,font=\normalfont\normalsize] at (-0.45,-1.95) {Có \textbf{#1} con. \textbf{#1} là số \SSocheck\ chẵn \quad\SSocheck\ lẻ};
  \end{tikzpicture}}
% Khung ghi nhớ số chẵn - số lẻ
\providecommand{\SSghinho}{}\renewcommand{\SSghinho}{%
  \begin{tikzpicture}[font=\normalfont]
    \draw[rounded corners=10pt,thick,draw=SSxanh!80,fill=SSxanh!7] (0,0) rectangle (16.4,5.9);
    \node[fill=SSxanh!80!SSxanhdam,text=white,rounded corners=4pt,font=\normalfont\bfseries,inner sep=3.5pt] at (2.0,5.9) {Ghi nhớ};
    \draw[SSxanh!60,thin,dashed] (8.2,1.6)--(8.2,5.3);
    % chẵn: 6 bọ rùa = 3 đôi
    \begin{scope}[shift={(2.95,4.4)}]\SSdoicore{6}{1}\end{scope}
    \node[align=center,text width=7.2cm,font=\normalfont\normalsize] at (4.1,2.15) {6 bạn bọ rùa ghép được 3 đôi, \textbf{không thừa} bạn nào.\\ Vậy \textbf{6 là số chẵn}.};
    % lẻ: 7 bọ rùa = 3 đôi và thừa 1
    \begin{scope}[shift={(10.6,4.4)}]\SSdoicore{7}{1}\end{scope}
    \node[align=center,text width=7.2cm,font=\normalfont\normalsize] at (12.3,2.15) {7 bạn bọ rùa ghép được 3 đôi, \textbf{thừa 1} bạn.\\ Vậy \textbf{7 là số lẻ}.};
    \node[draw=SSxanhdam,fill=SSnuoc!70,rounded corners=5pt,text width=7.4cm,align=center,inner sep=3pt,font=\normalfont\small] at (4.1,0.7) {Số chẵn có chữ số tận cùng là \textbf{0, 2, 4, 6, 8}.};
    \node[draw=SSdo!80!black,fill=SShong!35,rounded corners=5pt,text width=7.4cm,align=center,inner sep=3pt,font=\normalfont\small] at (12.3,0.7) {Số lẻ có chữ số tận cùng là \textbf{1, 3, 5, 7, 9}.};
  \end{tikzpicture}}

% Khung ghi nhớ "Mẹo vui" so sánh bằng cá sấu
\providecommand{\SSmeo}{}\renewcommand{\SSmeo}{%
  \begin{tikzpicture}[font=\normalfont]
    \draw[rounded corners=10pt,thick,draw=SSla!70!SSladam,fill=SSla!12] (0,0) rectangle (16.4,4.3);
    \draw[SSla!50!SSladam,thin,dashed] (8.2,1.1)--(8.2,3.2);
    \node[fill=SSla!70!SSladam,text=white,rounded corners=4pt,font=\normalfont\bfseries,inner sep=3.5pt] at (2.0,4.3) {Mẹo vui};
    \node[font=\normalfont\bfseries\normalsize] at (8.2,3.65) {Cá sấu luôn há miệng về phía số lớn hơn!};
    % ví dụ 1: 80 > 45 (cá sấu quay mặt về phía số lớn, bên trái)
    \begin{scope}[shift={(4.3,2.3)}]
      \node[SSthe=SSvang,circle,minimum size=1.15cm] at (-3.0,0) {80};
      \begin{scope}[shift={(0.2,-0.2)},xscale=-0.75,yscale=0.75]\SScasau\end{scope}
      \node[SSthe=SSxanh,circle,minimum size=1.15cm] at (3.0,0) {45};
      \node[font=\normalfont\bfseries\large] at (0,-1.4) {80 \textgreater\ 45};
    \end{scope}
    % ví dụ 2: 36 < 59 (cá sấu quay mặt về phía số lớn, bên phải)
    \begin{scope}[shift={(12.1,2.3)}]
      \node[SSthe=SSxanh,circle,minimum size=1.15cm] at (-3.0,0) {36};
      \begin{scope}[shift={(-0.2,-0.2)},scale=0.75]\SScasau\end{scope}
      \node[SSthe=SSvang,circle,minimum size=1.15cm] at (3.0,0) {59};
      \node[font=\normalfont\bfseries\large] at (0,-1.4) {36 \textless\ 59};
    \end{scope}
    \node[font=\normalfont\small] at (8.2,0.3) {Hai số bằng nhau thì dùng dấu =.};
  \end{tikzpicture}}

% Hàng so sánh trong thẻ: \SSsosanhhang{a/b, c/d, ...}  (4 thẻ mỗi hàng)
\providecommand{\SSsosanhluoi}{}\renewcommand{\SSsosanhluoi}[1]{%
  \begin{tikzpicture}[font=\normalfont\bfseries\Large]
    \foreach \a/\b [count=\k] in {#1}{%
      \pgfmathtruncatemacro{\cot}{mod(\k-1,4)}\pgfmathtruncatemacro{\hang}{div(\k-1,4)}\pgfmathtruncatemacro{\m}{mod(\cot,5)}\edef\mc{\SSmaule{\m}}%
      \begin{scope}[shift={(4.2*\cot,-1.2*\hang)}]
        \draw[rounded corners=8pt,thick,draw=\mc!70!black,fill=\mc!22] (0,-0.5) rectangle (3.9,0.5);
        \node at (0.8,0) {\a};
        \node[draw=\mc!70!black,thick,rounded corners=3pt,fill=white,minimum width=0.9cm,minimum height=0.75cm,inner sep=0pt] at (1.95,0) {};
        \node at (3.1,0) {\b};
      \end{scope}}
  \end{tikzpicture}}

%================ Bí ngô, cây thông, cây mùa thu ================
\providecommand{\SSbingo}{}\renewcommand{\SSbingo}{%
  \foreach \dx in {-0.55,0.55}{\shade[ball color=SScam!85!SSdo] (\dx,0) ellipse (0.52 and 0.74);}
  \shade[shading=radial,inner color=SSvang!55!SScam,outer color=SScam!90!SSdo] (0,0) ellipse (0.52 and 0.8);
  \foreach \x in {-0.3,0.3}{\draw[SSdo!55!SSnau,thin] (\x,0.72) .. controls (\x*1.5,0.3) and (\x*1.5,-0.3) .. (\x,-0.72);}
  \draw[SSdo!55!SSnau,thin] (-0.8,0.45) .. controls (-0.97,0.0) .. (-0.8,-0.45) (0.8,0.45) .. controls (0.97,0.0) .. (0.8,-0.45);
  \fill[SSnau!70!SSla,draw=SSnaudam,thin] (-0.1,0.76) .. controls (-0.12,0.95) and (0.0,1.05) .. (0.18,1.1) -- (0.24,0.98) .. controls (0.12,0.94) and (0.1,0.9) .. (0.1,0.77) -- cycle;
  \fill[SSla!80!SSladam,draw=SSladam,thin] (0.2,0.93) .. controls (0.45,1.2) and (0.75,1.05) .. (0.72,0.8) .. controls (0.5,0.85) and (0.3,0.85) .. (0.2,0.93);
  \draw[SSladam,thick,line cap=round] (0.2,0.98) .. controls (0.3,1.3) and (0.65,1.35) .. (0.7,1.12);}
% Khung bí ngô: \SShopbingo{tiêu đề}{t = bí ở góc trên trái, b = bí ở góc dưới phải}
\providecommand{\SShopbingo}{}\renewcommand{\SShopbingo}[2]{%
  \begin{tikzpicture}[font=\normalfont\bfseries]
    \draw[rounded corners=18pt,very thick,draw=SSdo!70!SSnau,fill=SSkem] (0,0) rectangle (7.8,4.6);
    \ifx#2t \node[text=SSdo!75!SSnau,font=\normalfont\bfseries\LARGE] at (4.3,3.95) {#1};
      \begin{scope}[shift={(0.35,4.25)},scale=0.85]\SSbingo\end{scope}
    \else \node[text=SSdo!75!SSnau,font=\normalfont\bfseries\LARGE] at (3.5,3.95) {#1};
      \begin{scope}[shift={(7.45,0.45)},scale=0.85]\SSbingo\end{scope}\fi
  \end{tikzpicture}}

% Cây thông Noel: \SScaythong{13 số}
\providecommand{\SScaythong}{}\renewcommand{\SScaythong}[1]{%
  \begin{tikzpicture}[font=\normalfont\bfseries\small]
    % quà dưới gốc
    \shade[left color=SSdo!70,right color=SSdo] (-2.7,0) rectangle (-1.6,0.75);\draw[SSdo!50!black,thin] (-2.7,0) rectangle (-1.6,0.75);
    \fill[SSvang] (-2.2,0) rectangle (-2.1,0.75) (-2.7,0.32) rectangle (-1.6,0.42);
    \shade[left color=SSxanh!70,right color=SSxanhdam] (1.5,0) rectangle (2.5,0.55);\draw[SSxanhdam!60!black,thin] (1.5,0) rectangle (2.5,0.55);
    \fill[white] (1.95,0) rectangle (2.05,0.55) (1.5,0.22) rectangle (2.5,0.32);
    \shade[left color=SStim!60,right color=SStim] (-1.25,0) rectangle (-0.65,0.4);\draw[SStim!50!black,thin] (-1.25,0) rectangle (-0.65,0.4);
    % thân cây
    \shade[left color=SSnau!70!SSvang,right color=SSnaudam] (-0.4,-0.05) rectangle (0.4,1.0);
    \draw[SSnaudam,thin] (-0.4,-0.05) rectangle (0.4,1.0);\draw[SSnaudam!70,very thin] (-0.15,0)--(-0.15,0.9) (0.12,0.05)--(0.12,0.95);
    % ba tầng tán lá (viền dưới răng cưa như cành thông)
    \foreach \y/\h/\w/\n in {0.9/2.7/3.0/7, 2.7/2.8/2.4/6, 4.6/2.4/1.7/5}{%
      \pgfmathsetmacro{\d}{2*\w/\n}%
      \shade[top color=SSla!75!white,bottom color=SSladam!95!SSla,rounded corners=2pt] (-\w,\y) -- (0,\y+\h) -- (\w,\y)
        \foreach \i in {1,...,\n}{ -- (\w-\i*\d+\d/2,\y-0.22) -- (\w-\i*\d,\y)} -- cycle;}
    % dây đèn vàng
    \draw[SSvang,thick,dash pattern=on 1pt off 3pt,line cap=round] (-2.5,1.1) .. controls (-1,0.7) and (1,0.7) .. (2.5,1.1);
    \draw[SSvang,thick,dash pattern=on 1pt off 3pt,line cap=round] (-1.9,2.9) .. controls (-0.8,2.55) and (0.8,2.55) .. (1.9,2.9);
    \draw[SSvang,thick,dash pattern=on 1pt off 3pt,line cap=round] (-1.35,4.8) .. controls (-0.5,4.5) and (0.5,4.5) .. (1.35,4.8);
    % ngôi sao
    \fill[SSvang,opacity=0.3] (0,7.25) circle (0.7);
    \node[star,star points=5,star point ratio=2.2,top color=SSvang!70!white,bottom color=SScam,draw=SScam!80!black,thick,minimum size=1.0cm,inner sep=0pt] at (0,7.25) {};
    % bóng có số (để tô màu)
    \foreach \x/\y [count=\k] in {0/6.2,-0.65/5.3,0.65/5.3,-0.8/4.0,0.8/4.0,-1.3/3.2,0/3.2,1.3/3.2,-0.9/2.0,0.9/2.0,-1.8/1.35,0/1.35,1.8/1.35}{%
      \node[circle,draw=SSxam!60!gray,very thick,fill=white,minimum size=0.82cm,inner sep=0pt] at (\x,\y) {\foreach \n [count=\j] in {#1}{\ifnum\j=\k\relax\n\fi}};}
  \end{tikzpicture}}

% Cây mùa thu: \SScaythu{10 số}
\providecommand{\SScaythu}{}\renewcommand{\SScaythu}[1]{%
  \begin{tikzpicture}[font=\normalfont\bfseries\small]
    % mặt đất
    \shade[top color=SSla!55!SSvang,bottom color=SSla!80!SSladam] (-3.5,0) .. controls (-1.8,0.7) and (1.8,0.7) .. (3.5,0) .. controls (1.8,-0.4) and (-1.8,-0.4) .. (-3.5,0);
    % thân và cành: lần 1 viền đậm, lần 2 màu gỗ
    \foreach \dw/\col in {3.2/SSnaudam,0/SSnau}{%
      \foreach \cd/\w in {{(0,0.1) .. controls (-0.12,1.0) and (0.1,1.8) .. (0,2.8)}/18,
                          {(0,2.1) .. controls (-0.8,2.6) and (-1.7,2.9) .. (-2.6,3.5)}/8,
                          {(0,2.5) .. controls (0.9,3.0) and (1.8,3.3) .. (2.7,3.7)}/8,
                          {(0,2.8) .. controls (-0.5,3.8) and (-0.9,4.6) .. (-1.2,5.3)}/8,
                          {(0,2.8) .. controls (0.5,3.8) and (0.8,4.6) .. (1.2,5.3)}/7,
                          {(0,1.5) .. controls (-0.8,1.8) and (-1.7,2.0) .. (-2.4,2.3)}/6,
                          {(0,1.3) .. controls (0.9,1.7) and (1.7,1.8) .. (2.5,2.2)}/6,
                          {(-0.35,3.35) .. controls (-0.3,4.2) and (-0.1,5.0) .. (0.15,5.7)}/5}{%
        \pgfmathsetmacro{\lw}{\w+\dw}\draw[\col,line width=\lw pt,line cap=round] \cd;}}
    \draw[SSnaudam!80,thin] (-0.25,0.5)--(-0.2,1.2) (0.15,0.8)--(0.1,1.7) (-0.1,1.6)--(-0.08,2.3);
    % lá rụng
    \foreach \x/\y/\a/\c in {-2.9/0.5/30/SSdo,2.5/0.35/-20/SScam,-1.5/0.15/60/SSvang,1.4/0.1/10/SSdo,3.0/1.6/70/SSvang,-3.1/1.3/-30/SScam}{%
      \begin{scope}[shift={(\x,\y)},rotate=\a]\fill[\c!85!SSnau,draw=SSnau!80,thin] (0,0.2) .. controls (0.2,0.15) and (0.2,-0.1) .. (0,-0.2) .. controls (-0.2,-0.1) and (-0.2,0.15) .. (0,0.2);\draw[SSnaudam,thin] (0,0.2)--(0,-0.26);\end{scope}}
    % lá có số: vị trí lá / đầu cành
    \foreach \x/\y/\tx/\ty [count=\k] in {-3.2/4.4/-2.6/3.5, -1.9/5.9/-1.2/5.3, 0.2/6.7/0.15/5.7, 1.9/5.9/1.2/5.3, 3.2/4.7/2.7/3.7, -3.2/2.7/-2.4/2.3, 3.2/2.6/2.5/2.2, -1.4/3.8/-1.1/2.8, 1.5/3.9/1.25/3.05, 0.0/4.7/0.0/4.2}{%
      \draw[SSnau,very thick,line cap=round] (\tx,\ty)--(\x,\y);
      \node[star,star points=5,star point ratio=1.45,draw=SSnau,very thick,fill=white,minimum size=1.4cm,inner sep=0pt,rounded corners=3pt] at (\x,\y) {\foreach \n [count=\j] in {#1}{\ifnum\j=\k\relax\n\fi}};}
  \end{tikzpicture}}

%================ Cá và xô, dãy nhà, liền trước - liền sau ================
\providecommand{\SSca}{}\renewcommand{\SSca}[2]{% màu, số (vẽ tại gốc, cá quay sang phải, dài khoảng 2,2cm)
  \begin{scope}[line join=round]
    \fill[#1!80!black,draw=#1!45!black,thick] (-0.7,0) -- (-1.15,0.42) .. controls (-1.0,0.15) and (-1.0,-0.15) .. (-1.15,-0.42) -- cycle;
    \fill[#1!75!black,draw=#1!45!black,thick] (-0.25,0.4) -- (-0.05,0.7) -- (0.35,0.42) -- cycle;
    \fill[#1!75!black,draw=#1!45!black,thick] (-0.1,-0.4) -- (0.0,-0.62) -- (0.25,-0.4) -- cycle;
    \path[top color=#1!85!white,bottom color=#1!30!white,draw=#1!45!black,thick] (0.98,0) .. controls (0.75,0.55) and (-0.35,0.6) .. (-0.72,0) .. controls (-0.35,-0.6) and (0.75,-0.55) .. (0.98,0);
    \draw[#1!50!black,thick] (0.52,0.37) .. controls (0.42,0.15) and (0.42,-0.15) .. (0.52,-0.36);
    \fill[white,draw=#1!45!black,thin] (0.76,0.13) circle (0.1);\fill[black] (0.79,0.13) circle (0.055);
    \node[font=\normalfont\bfseries\normalsize,fill=white,fill opacity=0.8,text opacity=1,rounded corners=3pt,inner sep=1.5pt] at (0.02,-0.02) {#2};
  \end{scope}}
% Xô đựng cá: \SSxo{màu}{chữ}
\providecommand{\SSxo}{}\renewcommand{\SSxo}[2]{%
  \draw[gray!70,line width=2pt] (-1.2,1.95) .. controls (-1.15,3.0) and (1.15,3.0) .. (1.2,1.95);
  \fill[black,opacity=0.12] (0,-0.05) ellipse (1.2 and 0.2);
  \shade[left color=#1!55!white,right color=#1!95!black] (-1.0,0.1) .. controls (-0.7,-0.2) and (0.7,-0.2) .. (1.0,0.1) -- (1.3,1.9) -- (-1.3,1.9) -- cycle;
  \draw[#1!50!black,thick] (-1.0,0.1) .. controls (-0.7,-0.2) and (0.7,-0.2) .. (1.0,0.1) -- (1.3,1.9) -- (-1.3,1.9) -- cycle;
  \fill[#1!40!black] (0,1.9) ellipse (1.3 and 0.27);\draw[#1!30!black,thick] (0,1.9) ellipse (1.3 and 0.27);
  \fill[#1!25!SSnuoc,opacity=0.9] (0,1.9) ellipse (1.15 and 0.2);
  \node[fill=white,draw=#1!50!black,rounded corners=4pt,minimum width=1.7cm,minimum height=0.6cm,inner sep=2pt,font=\normalfont\bfseries] at (0,0.95) {#2};}
% Bài nối cá với xô: \SSnoica{8 số}  (số chẵn: xô trên, số lẻ: xô dưới)
\providecommand{\SSnoica}{}\renewcommand{\SSnoica}[1]{%
  \begin{tikzpicture}
    % nước
    \shade[top color=SSnuoc!70!white,bottom color=SSxanh!55!SSnuoc,rounded corners=14pt] (-1.5,0.8) rectangle (2.4,-8.2);
    \draw[SSxanh!70,thick,rounded corners=14pt] (-1.5,0.8) rectangle (2.4,-8.2);
    \foreach \x/\y/\r in {-1.2/-0.3/0.09,-1.25/-2.6/0.07,-1.15/-4.4/0.11,-1.25/-6.2/0.08,-1.15/-7.8/0.1,2.05/-0.25/0.07}{\draw[white,thin] (\x,\y) circle (\r);}
    % cá
    \foreach \n [count=\k] in {#1}{%
      \pgfmathtruncatemacro{\m}{mod(\k-1,5)}\edef\mc{\SSmaule{\m}}%
      \begin{scope}[shift={(0,-1.0*\k+0.45)}]\SSca{\mc}{\n}\fill (1.75,0) circle (2.5pt);\end{scope}}
    % xô
    \begin{scope}[shift={(8.4,-3.0)}]\SSxo{SSxanhdam}{Số chẵn}\fill (-1.55,1.0) circle (2.5pt);\end{scope}
    \begin{scope}[shift={(8.4,-7.0)}]\SSxo{SSdo}{Số lẻ}\fill (-1.55,1.0) circle (2.5pt);\end{scope}
  \end{tikzpicture}}

% Ngôi nhà có biển số: \SSnha{màu tường}{màu mái}{số} (góc trái dưới tại gốc, rộng 2,1cm)
\providecommand{\SSnha}{}\renewcommand{\SSnha}[3]{%
  \fill[SSnau!70!SSnaudam] (1.4,1.4) rectangle (1.7,2.1);
  \shade[left color=#1!70!white,right color=#1] (0,0) rectangle (2.1,1.4);
  \draw[#1!55!black,thin] (0,0) rectangle (2.1,1.4);
  \shade[top color=#2!85!white,bottom color=#2!95!black,line join=round] (-0.2,1.4) -- (1.05,2.3) -- (2.3,1.4) -- cycle;
  \draw[#2!45!black,thin,line join=round] (-0.2,1.4) -- (1.05,2.3) -- (2.3,1.4) -- cycle;
  \fill[SSnau,draw=SSnaudam,thin] (0.22,0) rectangle (0.78,0.95);\fill[SSvang] (0.68,0.45) circle (0.03);
  \fill[SSnuoc,draw=SSnaudam,thin] (1.2,0.8) rectangle (1.85,1.25);\draw[SSnaudam,thin] (1.525,0.8)--(1.525,1.25) (1.2,1.025)--(1.85,1.025);
  \node[draw=SSnaudam,thick,fill=white,rounded corners=3pt,minimum width=0.95cm,minimum height=0.5cm,inner sep=0pt,font=\normalfont\bfseries\small] at (1.5,0.4) {#3};}
% Cây xanh nhỏ
\providecommand{\SScaynho}{}\renewcommand{\SScaynho}{%
  \fill[SSnau] (-0.06,0) rectangle (0.06,0.45);
  \shade[ball color=SSla!80!SSladam] (0,0.8) circle (0.38);}
% Hai dãy nhà: \SSpho{6 số hàng trên}{6 số hàng dưới} (để trống ô bằng {}, ví dụ {2,4,,8,,12})
\providecommand{\SSpho}{}\renewcommand{\SSpho}[2]{%
  \begin{tikzpicture}[font=\normalfont]
    \foreach \hang/\ds/\tenhang/\mauhang in {0/{#1}/{Dãy nhà bên trái đường}/SSxanh,-3.9/{#2}/{Dãy nhà bên phải đường}/SSdo}{%
      \begin{scope}[shift={(0,\hang)}]
        \shade[top color=SSla!45!white,bottom color=SSla!75!SSvang,rounded corners=8pt] (-0.2,-0.3) rectangle (16.6,0.3);
        \foreach \x in {1,...,5}{\begin{scope}[shift={(2.8*\x-0.1,0.05)},scale=0.8]\SScaynho\end{scope}}
        \foreach \so [count=\k] in \ds{%
          \pgfmathtruncatemacro{\m}{mod(\k-1,6)}%
          \edef\mt{\ifcase\m SSvang\or SSnuoc\or SShong\or SSla\or SSkem\or SStim\fi}%
          \edef\mm{\ifcase\m SSdo\or SSxanhdam\or SSnau\or SScam\or SSladam\or SSdo\fi}%
          \begin{scope}[shift={(2.8*\k-2.8+0.1,0.3)}]\SSnha{\mt}{\mm}{\so}\end{scope}}
        \node[anchor=south west,fill=\mauhang!80!black,text=white,rounded corners=3pt,inner sep=3pt,font=\normalfont\bfseries\small] at (-0.2,2.65) {\tenhang};
      \end{scope}}
  \end{tikzpicture}}

% Liền trước - số - liền sau: \SSlien{số}{liền trước}{liền sau}{1 = ô mẫu (đã điền, có khung)}
\providecommand{\SSlien}{}\renewcommand{\SSlien}[4]{%
  \begin{tikzpicture}[font=\normalfont\bfseries\large]
    \ifnum#4=1 \node[orange,font=\normalfont\bfseries\large] at (3.0,0.95) {Mẫu};\fi
    \node[SSo,minimum width=1.05cm,minimum height=0.9cm] (a) at (0,0) {\textcolor{red}{#2}};
    \node[SSthe=SSvang,circle,minimum size=1.2cm,inner sep=0pt] (b) at (2.0,0) {#1};
    \node[SSo,minimum width=1.05cm,minimum height=0.9cm] (c) at (4.0,0) {\textcolor{red}{#3}};
    \draw[-latex,very thick,SSxanhdam] (a.east)--(b.west);\draw[-latex,very thick,SSxanhdam] (b.east)--(c.west);
    \node[font=\normalfont\scriptsize\mdseries,anchor=north] at (0,-0.55) {liền trước};
    \node[font=\normalfont\scriptsize\mdseries,anchor=north] at (4.0,-0.55) {liền sau};
  \end{tikzpicture}}

%================ Ong tìm hoa (mê cung số lẻ) ================
\providecommand{\SSong}{}\renewcommand{\SSong}{% con ong quay sang phải, dài khoảng 1,5cm
  \begin{scope}[line join=round]
    \fill[SSnuoc!70!white,draw=SSxanh!70,thin,opacity=0.9,rotate around={25:(-0.05,0.2)}] (-0.05,0.35) ellipse (0.2 and 0.38);
    \fill[SSnuoc!70!white,draw=SSxanh!70,thin,opacity=0.9,rotate around={-15:(0.2,0.2)}] (0.2,0.35) ellipse (0.18 and 0.34);
    \draw[black,thick,line cap=round] (0.62,0.1)--(0.8,0.3) (0.66,0.05)--(0.88,0.12);
    \fill[black] (0.8,0.3) circle (0.03) (0.88,0.12) circle (0.03);
    \shade[left color=SSvang,right color=SSvang!70!SScam] (0,0) ellipse (0.5 and 0.33);
    \begin{scope}\clip (0,0) ellipse (0.5 and 0.33);\fill[black] (-0.3,-0.4) rectangle (-0.1,0.4) (0.05,-0.4) rectangle (0.25,0.4);\end{scope}
    \draw[black!70,thin] (0,0) ellipse (0.5 and 0.33);
    \fill[black] (-0.5,0) -- (-0.68,0.03) -- (-0.68,-0.03) -- cycle;
    \shade[ball color=SSvang!60!SSnau] (0.5,0.03) circle (0.22);
    \fill[white] (0.57,0.1) circle (0.07);\fill[black] (0.59,0.1) circle (0.035);
    \draw[black,thin,line cap=round] (0.55,-0.06) arc (-60:-120:0.12);
  \end{scope}}
\providecommand{\SShoa}{}\renewcommand{\SShoa}{% bông hoa có thân, tâm tại gốc
  \draw[SSladam,line width=2.4pt,line cap=round] (0,-0.2)--(0,-1.0);
  \fill[SSla,draw=SSladam,thin] (0,-0.75) .. controls (0.3,-0.55) and (0.5,-0.7) .. (0.45,-0.9) .. controls (0.2,-0.95) and (0.05,-0.85) .. (0,-0.75);
  \foreach \a in {0,45,...,315}{\shade[ball color=SShong!90!SSdo,rotate=\a] (0,0.38) ellipse (0.17 and 0.26);}
  \shade[ball color=SSvang] (0,0) circle (0.22);\fill[SScam!70!SSnau] (-0.06,0.05) circle (0.02) (0.07,0.0) circle (0.02) (0,-0.09) circle (0.02);}
% Mê cung: \SSmecung{{6 số},...} gồm 5 hàng, mỗi hàng 6 số
\providecommand{\SSmecung}{}\renewcommand{\SSmecung}[1]{%
  \begin{tikzpicture}[font=\normalfont\bfseries\Large]
    \foreach \hang [count=\i] in {#1}{\foreach \so [count=\j] in \hang{%
      \draw[rounded corners=8pt,thick,draw=SScam!80!SSvang,fill=SSvang!25] (1.85*\j-0.85,-1.2*\i-0.5) rectangle ++(1.7,1.0);
      \node at (1.85*\j,-1.2*\i) {\so};}}
    \begin{scope}[shift={(-1.0,-1.2)},scale=1.3]\SSong\end{scope}
    \node[font=\normalfont\bfseries\small,text=SScam!60!SSnau] at (-1.0,-0.4) {Ong thợ};
    \draw[-latex,very thick,SScam!80!SSnau] (-0.15,-1.2)--(0.35,-1.2);
    \draw[-latex,very thick,SScam!80!SSnau] (11.3,-6.0)--(11.95,-6.0);
    \begin{scope}[shift={(13.0,-5.6)},scale=1.4]\SShoa\end{scope}
    \node[font=\normalfont\bfseries\small,text=SSdo!70!black] at (13.0,-4.65) {Hoa mật};
  \end{tikzpicture}}

% Ba thẻ số: \SSlabai{a}{b}{c}
\providecommand{\SSlabai}{}\renewcommand{\SSlabai}[3]{%
  \begin{tikzpicture}[font=\normalfont\bfseries,scale=0.85,transform shape]
    \foreach \d/\x/\a/\c in {#1/0/-8/SSdo,#2/2.7/0/SSxanhdam,#3/5.4/8/SSladam}{%
      \begin{scope}[shift={(\x,0)},rotate=\a]
        \fill[black,opacity=0.12,rounded corners=7pt] (-0.82,-1.18) rectangle (0.92,1.3);
        \draw[rounded corners=7pt,very thick,draw=\c,fill=white] (-0.85,-1.25) rectangle (0.85,1.25);
        \node[text=\c,font=\normalfont\bfseries\fontsize{40}{44}\selectfont] at (0,0) {\d};
        \node[text=\c,font=\normalfont\bfseries\small] at (-0.55,0.95) {\d};
        \node[text=\c,font=\normalfont\bfseries\small,rotate=180] at (0.55,-0.95) {\d};
      \end{scope}}
  \end{tikzpicture}}

%%%%%%%%%%%%%%%%%%%%%%%%%%%%%%%%%%%%%%%%%%%%%%%%%%%%%%%%%%
%                       BUỔI 1
%%%%%%%%%%%%%%%%%%%%%%%%%%%%%%%%%%%%%%%%%%%%%%%%%%%%%%%%%%
%\setcounter{section}{3} % không chỉnh sửa
\hktuan{1}{17}
\par\medskip % xuống dòng: tiêu đề không nằm chung dòng với khung HK1 - Tuần 17
{\centering
        {\color{mauchinh!120}\bfseries\Large PHIẾU BÀI TẬP TOÁN TƯ DUY CẤP ĐỘ A}\\[3pt]
        \SStieude{CHỦ ĐỀ: SO SÁNH CÁC SỐ TRONG PHẠM VI 100 (BUỔI 1)}\par}
\vspace{6pt}
\TTthongtin

\begin{bt}
    Điền dấu \textgreater, \textless, = thích hợp vào ô trống.
    \par\nopagebreak\vspace{4pt}\nopagebreak{\centering\SSmeo\par}\vspace{8pt}
    {\centering\SSsosanhluoi{10/20,35/48,89/98,12/14,100/100,60/40,29/28,90/91,100/10,50/50,64/46,36/36,100/99,35/45,47/37,27/30,48/49,36/46,68/69,52/62,69/79,55/52,38/37,83/82}\par}
\end{bt}

\begin{bt}
    Nhìn hình, viết số rồi điền dấu \textgreater, \textless, = vào ô trống.
    \par\nopagebreak\vspace{4pt}\nopagebreak
    \noindent\SSsosanh{2}{3}{3}{2}\hfill\SSsosanh{4}{5}{3}{8}\par\medskip
    \noindent\SSsosanh{2}{7}{2}{7}\hfill\SSsosanh{6}{1}{5}{9}\par
\end{bt}

\begin{bt}
    Khoanh tròn vào số lớn nhất trong mỗi hàng.
    \par\nopagebreak\vspace{22pt}\nopagebreak
    \noindent\begin{minipage}[t]{0.49\linewidth}\SShangsomau{12,62,9,10,18}{2}\end{minipage}\hfill
    \begin{minipage}[t]{0.49\linewidth}\SShangso{88,46,9,10,92}\end{minipage}\par\smallskip
    \noindent\begin{minipage}[t]{0.49\linewidth}\SShangso{4,60,27,35,48}\end{minipage}\hfill
    \begin{minipage}[t]{0.49\linewidth}\SShangso{20,78,46,24,16}\end{minipage}\par\smallskip
    \noindent\begin{minipage}[t]{0.49\linewidth}\SShangso{100,98,4,29,50}\end{minipage}\hfill
    \begin{minipage}[t]{0.49\linewidth}\SShangso{58,36,99,5,30}\end{minipage}\par
\end{bt}

\begin{bt}
    Khoanh tròn vào số bé nhất trong mỗi hàng.
    \par\nopagebreak\vspace{6pt}\nopagebreak
    \noindent\begin{minipage}[t]{0.49\linewidth}\SShangso{25,7,80,30,100}\end{minipage}\hfill
    \begin{minipage}[t]{0.49\linewidth}\SShangso{90,10,28,58,45}\end{minipage}\par\smallskip
    \noindent\begin{minipage}[t]{0.49\linewidth}\SShangso{54,45,23,79,16}\end{minipage}\hfill
    \begin{minipage}[t]{0.49\linewidth}\SShangso{12,9,99,76,26}\end{minipage}\par\smallskip
    \noindent\begin{minipage}[t]{0.49\linewidth}\SShangso{38,40,99,27,89}\end{minipage}\hfill
    \begin{minipage}[t]{0.49\linewidth}\SShangso{3,76,23,14,64}\end{minipage}\par
\end{bt}

\begin{bt}
    Sắp xếp các số trên mình mỗi chú sâu đo vào các ô trống theo thứ tự từ bé đến lớn.
    \par\nopagebreak\vspace{4pt}\nopagebreak{\centering
        \SSsaudohang{SSla}{23,35,21,12,27}{bé nhất}{lớn nhất}{1}\par\vspace{2pt}
        \SSsaudohang{SSla}{34,50,62,42,67}{bé nhất}{lớn nhất}{1}\par\vspace{2pt}
        \SSsaudohang{SSla}{26,71,73,56,33}{bé nhất}{lớn nhất}{1}\par}
\end{bt}

\begin{bt}
    Sắp xếp các số trên mình mỗi chú sâu đo vào các ô trống theo thứ tự từ lớn đến bé.
    \par\nopagebreak\vspace{4pt}\nopagebreak{\centering
        \SSsaudohang{SSxanh}{21,31,18,28,15}{lớn nhất}{bé nhất}{0}\par\vspace{2pt}
        \SSsaudohang{SSxanh}{33,23,56,42,37}{lớn nhất}{bé nhất}{0}\par\vspace{2pt}
        \SSsaudohang{SSxanh}{62,31,23,57,39}{lớn nhất}{bé nhất}{0}\par}
\end{bt}

\begin{bt}
    Bác Hoa bày bán 4 loại trái cây trong các thùng gỗ như hình vẽ. Em hãy trả lời các câu hỏi sau.
    \par\nopagebreak\vspace{4pt}\nopagebreak{\centering
        \SSthung{\SStaoicon}{Táo}{64}\hfill\SSthung{\SScamicon}{Cam}{46}\hfill\SSthung{\SSleicon}{Lê}{69}\hfill\SSthung{\SSdauicon}{Dâu}{60}\par}
    \begin{listEX}
        \item Loại quả nào có nhiều quả nhất? \TTcham[4cm]
        \item Loại quả nào có ít quả nhất? \TTcham[4cm]
        \item Viết số quả của các loại trái cây theo thứ tự từ bé đến lớn: \TTcham[5.5cm]
        \item Điền dấu \textgreater, \textless, =: \hspace{0.5cm} 64 \SSo\ 46 \hspace{0.8cm} 60 \SSo\ 69 \hspace{0.8cm} 46 \SSo\ 60
    \end{listEX}
\end{bt}

\begin{bt}
    Tìm số hoặc chữ số thích hợp điền vào ô trống.
    \begin{listEX}
        \item 45 \textless\ \SSomau{46}\ \textless\ 47 
        \item \SShaicot{69 \textgreater\ \SSo\ \textgreater\ 67}{52 \textless\ \SSo\ \textless\ 54}
        \item \SShaicot{7\SSo\ \textgreater\ 78}{\SSo 5 \textgreater\ 85}
        \item Các chữ số có thể điền vào ô trống để 4\SSo\ \textless\ 42 là: \TTcham[3cm]
        \item Viết các số lớn hơn 47 và bé hơn 53: \TTcham[5cm]
    \end{listEX}
\end{bt}

\begin{bt}
    Ba bạn Mai, Hà, Nam thi nhảy dây. Số lần nhảy của Mai là 74, của Hà là 47, của Nam là 77. Em hãy viết tên bạn vào bục nhận giải (bạn nhảy được nhiều lần hơn thì xếp hạng cao hơn).
    \par\nopagebreak\vspace{4pt}\nopagebreak{\centering\SSbuc\par}
\end{bt}

\begin{bt}
    Đố vui: Em là thám tử số học. Hãy tìm câu trả lời!
    \begin{listEX}
        \item Lan có 38 viên bi, Bình có 40 viên bi. Nam có số bi nhiều hơn Lan nhưng ít hơn Bình. Hỏi Nam có bao nhiêu viên bi?\\[4pt]
        Nam có: \TTcham[2.5cm] viên bi.
        \item Số lớn nhất có hai chữ số khác nhau là: \TTcham[2cm]
        \item Số bé nhất có hai chữ số là: \TTcham[2cm]
        \item Số có hai chữ số, lớn hơn 54, bé hơn 60 và có chữ số hàng đơn vị là 7 là số: \TTcham[2cm]
        \item Có bao nhiêu số lớn hơn 45 và bé hơn 50? \TTcham[2cm]
    \end{listEX}
\end{bt}

%%%%%%%%%%%%%%%%%%%%%%%%%%%%%%%%%%%%%%%%%%%%%%%%%%%%%%%%%%
%                       BUỔI 2
%%%%%%%%%%%%%%%%%%%%%%%%%%%%%%%%%%%%%%%%%%%%%%%%%%%%%%%%%%
\newpage
\TTdatlaibai
{\centering
        {\color{mauchinh!120}\bfseries\Large PHIẾU BÀI TẬP TOÁN TƯ DUY CẤP ĐỘ A}\\[3pt]
        \SStieude{CHỦ ĐỀ: SO SÁNH CÁC SỐ TRONG PHẠM VI 100 (BUỔI 2)}\par}
\vspace{6pt}
\TTthongtin

\begin{bt}
    Khoanh vào từng đôi bọ rùa rồi cho biết số bọ rùa là số chẵn hay số lẻ (đánh dấu vào ô thích hợp).
    \par\nopagebreak\vspace{4pt}\nopagebreak{\centering\SSghinho\par}\vspace{8pt}
    \noindent\SSdoibai{8}\hfill\SSdoibai{11}\par\smallskip
    \noindent\SSdoibai{12}\hfill\SSdoibai{9}\par
\end{bt}

\begin{bt}
    Khoanh tròn vào số chẵn.
    \par\nopagebreak\vspace{6pt}\nopagebreak{\centering
        \SSbonghang{4,5,6,8,9,10,14,17,11,52,36,15,19,28,71,34,12,1,58,18,33,27,10,25}\par}
\end{bt}

\begin{bt}
    Khoanh tròn vào số lẻ.
    \par\nopagebreak\vspace{6pt}\nopagebreak{\centering
        \SShuyhieuhang{1,32,5,8,4,16,77,31,90,45,6,37,58,31,9,10,36,55,89,100,2,35,76,44}\par}
\end{bt}

\begin{bt}
    Cho các số: 2; 7; 10; 16; 28; 33; 81; 90; 57. Em hãy sắp xếp các số trên vào ô số chẵn, số lẻ.
    \par\nopagebreak\vspace{4pt}\nopagebreak{\centering
    \begin{tikzpicture}
        \foreach \x [count=\k] in {2,7,10,16,28,33,81,90,57}{%
          \pgfmathtruncatemacro{\m}{mod(\k-1,5)}\edef\mc{\SSmaule{\m}}%
          \node[SSthe=\mc] at (1.45*\k,0) {\x};}
    \end{tikzpicture}\par\vspace{8pt}
    \SShopbingo{Số chẵn}{t}\hfill\SShopbingo{Số lẻ}{b}\par}
\end{bt}

\begin{bt}
    Tô màu đỏ vào số chẵn, màu xanh dương vào số lẻ.
    \par\nopagebreak\vspace{4pt}\nopagebreak{\centering
        \SScaythong{75,26,54,37,41,97,70,82,64,93,25,15,38}\hfill\SScaythu{74,22,54,41,89,68,76,25,13,37}\par}
\end{bt}

\newpage
\begin{bt}
    Chú mèo con đi câu được 8 chú cá. Em hãy nối mỗi chú cá với xô thích hợp (xô số chẵn hoặc xô số lẻ).
    \par\nopagebreak\vspace{4pt}\nopagebreak{\centering\SSnoica{14,27,50,63,8,35,72,91}\par}
\end{bt}

\begin{bt}
    Ở phố Hoa Sen, các nhà bên trái đường được đánh số chẵn, các nhà bên phải đường được đánh số lẻ. Em hãy viết số nhà còn thiếu vào biển số.
    \par\nopagebreak\vspace{4pt}\nopagebreak{\centering\SSpho{2,4,,8,,12}{1,,5,,9,}\par}
\end{bt}

\begin{bt}
    Viết số liền trước và số liền sau của mỗi số vào ô trống.
    \par\nopagebreak\vspace{4pt}\nopagebreak{\centering
        \SSlien{29}{28}{30}{1}\hfill\SSlien{54}{}{}{0}\hfill\SSlien{80}{}{}{0}\par\vspace{10pt}
        \SSlien{99}{}{}{0}\hfill\SSlien{36}{}{}{0}\hfill\SSlien{61}{}{}{0}\par}
\end{bt}

\newpage
\begin{bt}
    Giúp ong thợ bay về bông hoa mật: ong chỉ được bay qua các ô ghi số lẻ, mỗi lần bay sang một ô bên cạnh (trái, phải, lên hoặc xuống). Em hãy tô màu đường đi của ong.
    \par\nopagebreak\vspace{4pt}\nopagebreak{\centering
        \SSmecung{{81,47,66,51,84,75},{33,70,91,89,23,67},{49,28,21,78,98,15},{87,60,93,88,63,31},{11,77,19,69,14,35}}\par}
\end{bt}

\begin{bt}
    Ba tấm thẻ số 3, 5, 8 được xếp thành các số có hai chữ số khác nhau. Em hãy viết số theo yêu cầu.
    \par\nopagebreak\vspace{4pt}\nopagebreak{\centering\SSlabai{3}{5}{8}\par}
    \begin{listEX}
        \item \SShaicot{Số chẵn lớn nhất là: \TTcham[2cm]}{Số chẵn bé nhất là: \TTcham[2cm]}
        \item \SShaicot{Số lẻ lớn nhất là: \TTcham[2cm]}{Số lẻ bé nhất là: \TTcham[2cm]}
    \end{listEX}
\end{bt}

\begin{bt}
    Đố vui: Mình là số nào?
    \begin{listEX}
        \item Mình là số chẵn, lớn hơn 58 và bé hơn 61. Mình là số: \TTcham[2cm]
        \item Mình là số lẻ có hai chữ số, chữ số hàng chục là 8 và mình lớn hơn 87. Mình là số: \TTcham[2cm]
        \item Mình là số chẵn lớn nhất có hai chữ số. Mình là số: \TTcham[2cm]
        \item Lớp 1A có 27 bạn xếp thành từng đôi để tập múa. Hỏi có bao nhiêu bạn chưa có bạn cùng đôi?\\[4pt]
        Trả lời: \TTcham[3cm] bạn chưa có bạn cùng đôi.
    \end{listEX}
\end{bt}

\begin{bt}
    (Không bắt buộc) Điền số thích hợp vào chỗ trống.
    \begin{listEX}
        \item \SShaicot{Số liền trước của 12 là: \TTcham[1.8cm]}{Số liền sau của 68 là: \TTcham[1.8cm]}
        \item Các số nằm giữa 8 và 14 là: \TTcham[5cm]
        \item \SShaicot{Số chẵn liền sau của 10 là: \TTcham[1.8cm]}{Số lẻ liền trước của 28 là: \TTcham[1.8cm]}
        \item \SShaicot{Số lẻ liền sau của 35 là: \TTcham[1.8cm]}{Số chẵn ở giữa 5 và 10 là: \TTcham[2.2cm]}
        \item Số lẻ ở giữa 12 và 18 là: \TTcham[3cm]
    \end{listEX}
\end{bt}

\begin{center}
       \faGg\faGg\quad Hết \quad\faGg\faGg
\end{center}
\clearpage
\makeatletter\let\@textbottom\SSluutextbottom\makeatother
\end{document}
